\FloatBarrier
\section{Vintages and thresholds}\label{si:thresholds}

\paragraph{Vintage overlap, 1983--2015} The NCAM Table A2 vintage (Specification A) and the RAM Legacy v4.66 continuation of the same assessment family (stock COD2J3KL) differ over 1983--2015 by $+24.7\kt$ on average (Table A2 higher), by 26.8\,kt in mean absolute value and by 131.2\,kt at most, in 1984 (Fig.~\ref{fig:si_vintages}). The xteNCAM revision differs from Table A2 over the same years by 126.2\,kt RMSE, with Table A2 higher by 58.5\,kt on average. The 2015 origin is 298.65\,kt on Table A2, 277\,kt on RAM and 273\,kt on xteNCAM.

\paragraph{Vintage-shifted LRP} The LRP defined as the 1983--1989 mean is 884.6\,kt on the Table A2 vintage and 790.0\,kt on the RAM vintage, a shift of 94.6\,kt in the reference point from a revision of the same model family. Persistence from the 2015 origin under-predicts every outcome year 2016--2021 in both vintages: MAE 129.2\,kt and RMSE 133.3\,kt from the RAM origin of 277\,kt, MAE 107.5\,kt from the Table A2 origin. The outcome years agree across the RAM and xteNCAM vintages to within 1--26\,kt (mean absolute difference 16.8\,kt).

\paragraph{Threshold classification on one file} The xteNCAM file (1954--2024, 71 years) carries a ratio column that fixes $\Blim$ at 276\,kt (per-year implied values 250--289, mean 275.5). Classifying each year under $\Blim = 276\kt$ and under the 884.6-kt LRP gives 45 years at or above $\Blim$, 16 years at or above the LRP (1954--1969 only), and 29 disagreement years in two blocks: 1970--1992 without 1977--1979, and every outcome year 2016--2024. Under the 1980s-mean threshold the stock stands below the reference point for 55 consecutive years (1970--2024); under $\Blim$ it is below for 26 years (1977--1979 and 1993--2015) and above from 2016 onward. The late-1970s episode (ratios 0.69--0.95) is a critical-zone reading, on the current model, of years the contemporary fishery regarded as healthy. The checkpoints 2005, 2017 and 2024 read 26, 451 and 342\,kt. The October 2023 framework revision, which moved the stock's reported zone from critical to cautious without a change in the stock, is this experiment executed by the assessment authority. The scored RMSE comparison never uses the thresholds; every normative verdict layered on the forecasts, the Brier indicator included, inherits the dependence.

\begin{figure}[!htb]
\centering
\includegraphics[width=\textwidth]{figS2_vintages}
\caption{Three vintages of Northern cod SSB over 1983--2024: NCAM Table A2 (Specification A, 1983--2015), RAM Legacy v4.66 (the same assessment family continued to 2021) and xteNCAM Table 17 (Specification B). Horizontal lines mark the two Specification A LRP values (884.6\,kt on Table A2, 790.0\,kt on the RAM vintage) and $\Blim = 276\kt$. The outcome years 2016--2021 lie above the 2015 origin in every vintage and below 884.6\,kt in every vintage.}
\label{fig:si_vintages}
\end{figure}

\FloatBarrier
\section{The 2022--2023 observation suspension}\label{si:suspension}

The test conditions on a complete annual series. The stock's real 2022--2023 history violates that condition. Partial coverage by the fall research-vessel survey in 2021 and 2022 left the leading indicator of stock size uncomputable; no typical assessment update was possible in 2022 or 2023; the stewardship maximum was rolled over unchanged at 12{,}999 t; and the survey vessel that carries the series was condemned in February 2023. Table~\ref{tab:si_suspension} assembles the episode from the assessment of record and the public record.

\begin{table}[!htb]
\centering\small
\caption{The 2022--2023 observation suspension and its resolution (sources: Regular et al.\ 2025 and its Table 17 file; the 2024 reopening announcement; the 2025 and 2026 assessment coverage).}
\label{tab:si_suspension}
\begin{tabularx}{\textwidth}{@{}lL@{}}
\toprule
Period & Event \\
\midrule
2021--22 & Partial fall-survey coverage; leading indicator not computable \\
2022--23 & No typical assessment update; stewardship maximum held at 12{,}999 t \\
Feb 2023 & Survey vessel condemned \\
Mar 2024 & Revised-framework assessment: SSB at 1.2 \(B_{\mathrm{lim}}\) (95\% CI 0.7--2.1; 22\% critical-zone probability) \\
Jun 2024 & Commercial fishery reopened at 18{,}000 t after 32 years \\
2025 & TAC 38{,}000 t; reference point revised after being deemed overestimated \\
Apr 2026 & Healthy-zone report: SSB $\approx$ 540 kt (CI 420--700); 70\% probability above the upper stock reference \\
\bottomrule
\end{tabularx}
\end{table}

Read against a review-timing rule under which a review must land no later than the worst-case loss of the information it consumes, the episode is that rule's boundary case in the wild: the information loss ran past the review cadence, and the removals rule froze for two cycles while the observation map was dark. Abaee (2026, Periodic review as sampled governance) states the formal identity; the 2022--2023 suspension is its empirical instance, and it marks the practical limit of any scoring exercise that presumes the series exists.

\FloatBarrier
\section{Reproducibility and verification record}\label{si:repro}

\paragraph{Determinism} All computations are deterministic: within a fixed interpreter and library stack, repeated executions reproduce every result file byte for byte. An independent re-execution verified all 29 recorded output checksums and regenerated all 30 result files byte-identically; one archived file carries no pinned checksum and is covered by the regeneration comparison. Across environments (Python 3.12, numpy 2.1.3, scipy 1.14.1 against the recorded original environment), the intervention runners and the OLS-based prediction runners regenerate their outputs byte for byte, while the four optimiser-based forecast runners (L-BFGS-B fits) reproduce every scored row at printed precision with one exception: the M1b rolling row on Specification B, whose depensatory optimum is environment-sensitive at the $\pm17\kt$ level ($h=1$: 151.6 against 153.2\,kt; $h=5$: 445.5 against 462.5\,kt). That sensitivity is the identification fragility of Section~\ref{si:params} and is immaterial to the verdict, which persistence wins at both horizons by margins far larger than the instability.

\paragraph{Third environment (preceding edition)} For the preceding single-system edition the four forecast runners (\texttt{run\_ladder.py}, \texttt{run\_xte.py}, \texttt{run\_capelin\_regime.py}, \texttt{run\_capelin\_index.py}), the uncertainty script and the simulation outputs were re-executed or re-read in a third environment (Python 3.13.16, numpy 2.3.5, scipy 1.17.1). Every result file, including the M1b rows, was reproduced at printed precision, and the persistence baselines were recomputed from the series on the archived origin sets and matched the archived origin-matched values (98.05, 264.72, 84.43, 299.98\,kt). The paired block bootstrap was re-run from the archived per-origin files with an explicitly shared index draw and reproduced every interval of Table~\ref{tab:si_dm}. The period decomposition, the rolling parameter record, the floor census and the catch-timing sensitivity were computed in the same environment from the archived files and the unmodified runner functions; their scripts are archived under \texttt{paperE1\_cod\_forecast\_ladder\_v69\_src/work/}. The present edition re-reads the archived result files and recomputes the two baselines on the archived origin sets (S23); no forecast runner was re-executed and no archived cod score changed.

\paragraph{Verification script} \texttt{paper11\_forecasting\_\allowbreak{}baselines\_v67\_\allowbreak{}verification.py} (standard library only) regenerates every number of main-text Section 3.7, of S10 and of Sections~\ref{si:thresholds}--\ref{si:suspension} from the archived data files (NCAM Table A2, RAM v4.66 time series, xteNCAM Table 17, the capelin index) and asserts the cross-checks, including the presence in the manuscript source of the pinned values and the suspension facts. All checks pass on the main text and this document taken together.

\paragraph{Scripts} (\texttt{SRC} denotes the source folder \texttt{paper11\_forecasting\_\allowbreak{}baselines\_v67\_src} of the repository.)
{\footnotesize\begin{verbatim}
python3 wave_e_cod/src/run_ladder.py
python3 wave_e_cod/src/run_xte.py
python3 wave_e_cod/src/run_capelin_regime.py
python3 wave_e_cod/src/run_capelin_index.py
python3 "campaign_e1_dm_uncertainty_regime.py"
python3 tools/sim_retention_power.py
python3 SRC/work/cod/paired_bootstrap.py
python3 SRC/work/cod/post1995.py
python3 SRC/work/cod/catch_timing_sensitivity.py
python3 SRC/work/cod/make_figures.py
\end{verbatim}}

\FloatBarrier
\section{Catch-timing convention: a sensitivity outside the archived protocol}\label{si:timing}

In \texttt{run\_ladder.py} the estimator pairs the increment $S_{j+1}-S_j$ with the catch $C_j$ of the same year, whereas the forecast path at origin $t$ applies $C_{t+1},\dots,C_{t+h}$ to the steps $t\to t+1,\dots,t+h-1\to t+h$. The two conventions differ by one year. The estimation convention is the one under which SSB is a start-of-year quantity and $C_t$ the removals during year $t$; the forecast convention gives the catch-supplied modules the catch of the target year, which under that reading occurs after the predictand is realised. The offset was noticed during the re-execution for this revision. Table~\ref{tab:si_timing} re-runs M2--M4 with the estimation convention in the forecast path (persistence and M1/M1b are unaffected). Under the coarse regime the two conventions differ only at the 1990--1992 origins, where the regime boundaries fall; under annual landings and on Specification B they differ at every origin.

\begin{table}[!htb]
\centering\small
\caption{Rolling-origin RMSE (kt) of the catch-supplied modules under the archived forecast convention ($C_{t+1}$ applied to the step $t\to t+1$) and under the estimation convention ($C_t$), with the ratio to origin-matched persistence in parentheses.}
\label{tab:si_timing}
\begin{tabular}{@{}llcrrr@{}}
\toprule
Spec.\ (catch) & Module & $h$ & archived & estimation convention & persistence \\
\midrule
A annual & M2 & 1 & 160.4 (1.64) & 133.3 (1.36) & 98.0 \\
A annual & M2 & 5 & 393.8 (1.49) & 390.8 (1.48) & 264.7 \\
A annual & M3 & 1 & 153.6 (1.57) & 125.5 (1.28) & 98.0 \\
A annual & M3 & 5 & 351.5 (1.33) & 352.2 (1.33) & 264.7 \\
A annual & M4 & 1 & 206.3 (2.10) & 161.9 (1.65) & 98.0 \\
A annual & M4 & 5 & 486.4 (1.84) & 438.7 (1.66) & 264.7 \\
A coarse & M2 & 1 & 144.1 (1.47) & 113.4 (1.16) & 98.0 \\
A coarse & M2 & 5 & 398.2 (1.50) & 367.8 (1.39) & 264.7 \\
A coarse & M3 & 1 & 134.6 (1.37) & 105.2 (1.07) & 98.0 \\
A coarse & M3 & 5 & 366.4 (1.38) & 341.5 (1.29) & 264.7 \\
A coarse & M4 & 1 & 195.6 (1.99) & 138.9 (1.42) & 98.0 \\
A coarse & M4 & 5 & 488.2 (1.84) & 409.8 (1.55) & 264.7 \\
B & M2 & 1 & 166.0 (1.97) & 146.4 (1.73) & 84.4 \\
B & M2 & 5 & 1058.9 (3.53) & 981.8 (3.27) & 300.0 \\
B & M3 & 1 & 126.7 (1.50) & 94.9 (1.12) & 84.4 \\
B & M3 & 5 & 930.1 (3.10) & 832.5 (2.78) & 300.0 \\
B & M4 & 1 & 205.7 (2.44) & 130.3 (1.54) & 84.4 \\
B & M4 & 5 & 1030.7 (3.44) & 898.3 (2.99) & 300.0 \\
\bottomrule
\end{tabular}
\end{table}

Under the estimation convention every one-step deficit shrinks and no five-year deficit changes sign. M3 is 7\% above persistence at one year on Specification A (coarse) and 12\% on Specification B, still outside the 5\% band, and 29\% and 178\% above at five years. The post-1995 bands of Table 6 are unchanged on Specification A (the catch is constant at 5\,kt there under the coarse regime); on Specification B the post-1995 one-step M3 margin becomes $+10.4\kt$ with interval $[-20.2, +46.8]$ ($p = 0.53$), M2 and M4 remain significantly worse ($+62.0$, $[+24.6, +98.7]$; $+45.9$, $[+11.8, +86.4]$), and all five-year margins remain far outside the band. The verdict of the main text does not depend on the convention; the size of the one-step deficits does, and a future application of the ladder should use the estimation convention throughout.

\begin{figure}[!htb]
\centering
\includegraphics[width=\textwidth]{figS3_windows_annual}
\caption{Fixed-window forecast paths on Specification A with annual landings (companion to Fig.~\ref{fig:si_windows} in S4). (a) Collapse window: as under the coarse regime, every module forecasts continuation or growth. (b) Recovery window: with $K$ at its 5000-kt bound, production does not saturate and the catch-supplied modules over-predict more than under the coarse regime (M3 609\,kt RMSE against 220); M1b is again a flat line at its fitted $K$.}
\label{fig:si_windows_annual}
\label{fig:si_annual}
\end{figure}
